\chapter*{Résumé}
% ... Ajoutez ici le contenu du résumé ...
Ce rapport présente le Projet de Fin d'Année « Développement et conception d'une plateforme de gouvernance des risques algorithmiques ». Le projet couvre l'analyse des besoins, la conception et la réalisation opérationnelle d'une plateforme appelée PAGe Platform, destinée aux organisations publiques soucieuses de gouverner les risques algorithmiques.

La première partie définit les concepts et notions essentiels qui fondent toute l'approche : qu'entend-on par risque algorithmique, indicateur d'évaluation, mécanisme de mitigation (RMM), pilier de gouvernance et maturité institutionnelle. Elle établit le cadre conceptuel et méthodologique permettant aux organisations publiques d'évaluer et de suivre systématiquement les risques liés à l'utilisation d'algorithmes dans leurs processus décisionnels et opérationnels.

La deuxième partie décrit l'architecture générale et fonctionnelle de la plateforme, organisée autour de trois modules complémentaires : un premier module d'évaluation de la gravité et de l'exposition aux risques algorithmiques, un second module d'évaluation des capacités institutionnelles et des mécanismes de mitigation disponibles, et un troisième module d'analyse d'impact et de simulation des scénarios possibles. Elle expose en détail la logique d'agrégation des mesures élémentaires et le processus de calcul des scores de maturité institutionnelle.

La troisième partie rapporte et documente la phase de réalisation concrète du projet : modélisation et structuration des données, implémentation des modules fonctionnels, et validation rigoureuse des résultats. Elle présente les livrables produits (définitions formelles, tableaux d'agrégation, exemples numériques, règles de calcul détaillées) et décrit les principaux cas d'usage et scénarios supportés par la plateforme.

Le rapport conclut en fournissant des recommandations pratiques pour l'utilisation effective de la plateforme et l'interprétation correcte des résultats afin d'améliorer progressivement la gouvernance institutionnelle des risques algorithmiques.
\vspace{6cm}
\begin{quote}
\textbf{Mots-clés :} Gouvernance algorithmique, Indicateurs, Mitigation, Simulation, Tableau de bord.
\end{quote}